\textbf{Phần 1. Trắc nghiệm lựa chọn}
\setcounterpageref{ex}{0}

\begin{ex} %[3D11N60]
Điền dấu thích hợp: $45\,200 \dots 42\,500$.
\choice
{$<$}
{$=$}
{\True $>$}
{Không so sánh được}
\loigiai{
Hàng nghìn: $5 > 2$ nên $45200 > 42500$.\\
\textbf{Đáp án của câu hỏi là C}
}
\end{ex}

\begin{ex} %[3D11N60]
Điền dấu thích hợp: $89\,999 \dots 90\,000$.
\choice
{$>$}
{$=$}
{Không so sánh được}
{\True $<$}
\loigiai{
$89999 < 90000$.\\
\textbf{Đáp án của câu hỏi là D}
}
\end{ex}

\begin{ex} %[3D11H60]
Số nào sau đây lớn nhất?
\choice
{\True $65800$}
{$65430$}
{$65408$}
{$65084$}
\loigiai{
Số lớn nhất là $65800$.\\
\textbf{Đáp án của câu hỏi là A}
}
\end{ex}

\begin{ex} %[3D11H60]
Số nào sau đây bé nhất?
\choice
{$78650$}
{$78065$}
{\True $70865$}
{$76850$}
\loigiai{
Số bé nhất là $70865$.\\
\textbf{Đáp án của câu hỏi là C}
}
\end{ex}

\begin{ex} %[3D11H60]
Sắp xếp các số $45600, 46500, 40650, 45060$ theo thứ tự từ bé đến lớn:
\choice
{$46500 < 45600 < 45060 < 40650$}
{$40650 < 45600 < 45060 < 46500$}
{$45060 < 40650 < 45600 < 46500$}
{\True $40650 < 45060 < 45600 < 46500$}
\loigiai{
Thứ tự tăng dần: $40650 < 45060 < 45600 < 46500$.\\
\textbf{Đáp án của câu hỏi là D}
}
\end{ex}

\begin{ex} %[3D11H60]
Điền chữ số thích hợp vào ô trống để $68\,54\dots > 68\,548$:
\choice
{\True $9$}
{$7$}
{$8$}
{$6$}
\loigiai{
Chữ số hàng đơn vị phải là 9.\\
\textbf{Đáp án của câu hỏi là A}
}
\end{ex}

\begin{ex} %[3D11V60]
Trong các số $78\,450; 87\,450; 78\,540; 87\,540$, số lớn nhất là
\choice
{$87450$}
{\True $87540$}
{$78540$}
{$78450$}
\loigiai{
Số lớn nhất là $87540$.\\
\textbf{Đáp án của câu hỏi là B}
}
\end{ex}

\begin{ex} %[3D11V60]
Tìm số tự nhiên $x$ tròn chục nghìn biết: $35\,000 < x < 55\,000$.
\choice
{\True $40000$ (hoặc $50000$)}
{$30000$}
{$60000$}
{$20000$}
\loigiai{
Các số tròn chục nghìn thỏa mãn là $40000$ hoặc $50000$.\\
\textbf{Đáp án của câu hỏi là A}
}
\end{ex}

\begin{ex} %[3D11V60]
Một cái ti vi giá $75\,000$ đồng (hoặc $7\,500\,000$), một cái quạt giá $45\,000$ đồng. Thiết bị nào có giá đắt hơn?
\choice
{Cái quạt}
{Bằng nhau}
{Không so sánh được}
{\True Cái ti vi}
\loigiai{
Vì $75000 > 45000$ nên ti vi đắt hơn.\\
\textbf{Đáp án của câu hỏi là D}
}
\end{ex}

\begin{ex} %[3D11V60]
Tìm chữ số $x$ để: $4x500 < 41500$.
\choice
{$1$}
{\True $0$}
{$2$}
{$3$}
\loigiai{
$x < 1 \Rightarrow x = 0$.\\
\textbf{Đáp án của câu hỏi là B}
}
\end{ex}

\begin{ex} %[3D11V60]
Dãy số nào sau đây xếp theo thứ tự giảm dần?
\choice
{$12000; 13000; 14000; 15000$}
{$32000; 35000; 34000; 36000$}
{\True $95000; 94000; 89000; 87000$}
{$50000; 60000; 70000; 80000$}
\loigiai{
Dãy giảm dần: $95000 > 94000 > 89000 > 87000$.\\
\textbf{Đáp án của câu hỏi là C}
}
\end{ex}

\begin{ex} %[3D11C60]
Tìm số tròn trăm $x$ có 5 chữ số biết: $34\,850 < x < 35\,050$.
\choice
{$34800$}
{\True $34900$ (hoặc $35000$)}
{$35100$}
{$34700$}
\loigiai{
Số tròn trăm thỏa mãn là $34900$ hoặc $35000$.\\
\textbf{Đáp án của câu hỏi là B}
}
\end{ex}

\textbf{Phần 2. Trắc nghiệm đúng sai}
\setcounterpageref{ex}{0}

\begin{ex} %[3D11H60]
Xét tính đúng sai khi so sánh các số:
\choiceTF
{\True $54320 < 54325$.}
{\True $100000 > 99999$.}
{$60000 + 4000 > 64005$.}
{\True $60000 + 4000 < 64005$ (vì $64000 < 64005$).}
\loigiai{
\begin{itemize}
    \item $54320 < 54325$: Đúng.
    \item $100000 > 99999$: Đúng.
    \item $64000 < 64005$. Ý 3 Sai, Ý 4 Đúng.
\end{itemize}
\begin{itemize}
    \item a) \textbf{Đúng}.
    \item b) \textbf{Đúng}.
    \item c) \textbf{Sai}.
    \item d) \textbf{Đúng}.
\end{itemize}
}
\end{ex}

\begin{ex} %[3D11V60]
Cho 4 số: $45\,600; 46\,500; 45\,060; 46\,050$. Xét tính đúng sai:
\choiceTF
{\True Số lớn nhất là $46\,500$.}
{\True Số bé nhất là $45\,060$.}
{Có 3 số lớn hơn $46\,000$.}
{\True Có 2 số lớn hơn $46\,000$ (là $46\,050$ và $46\,500$).}
\loigiai{
\begin{itemize}
    \item Lớn nhất: $46500$. Đúng.
    \item Bé nhất: $45060$. Đúng.
    \item Lớn hơn 46000 có 2 số. Ý 3 Sai, Ý 4 Đúng.
\end{itemize}
\begin{itemize}
    \item a) \textbf{Đúng}.
    \item b) \textbf{Đúng}.
    \item c) \textbf{Sai}.
    \item d) \textbf{Đúng}.
\end{itemize}
}
\end{ex}

\textbf{Phần 3. Trả lời ngắn}
\setcounterpageref{ex}{0}

\begin{ex} %[3D11H60]
Tìm số lớn nhất trong các số: $78\,450; 78\,540; 78\,405; 78\,504$.
\shortans{78540}
\loigiai{
\textbf{Đáp số:} $78540$\\
$78540$.
}
\end{ex}

\begin{ex} %[3D11H60]
Tìm số bé nhất trong các số: $91\,200; 90\,120; 92\,010; 91\,020$.
\shortans{90120}
\loigiai{
\textbf{Đáp số:} $90120$\\
$90120$.
}
\end{ex}

\begin{ex} %[3D11V60]
Tìm chữ số $x$ để: $52\,x40 < 52\,140$.
\shortans{0}
\loigiai{
\textbf{Đáp số:} $0$\\
$x = 0$.
}
\end{ex}

\begin{ex} %[3D11V60]
Tìm số tròn chục nghìn $x$ biết: $65\,000 < x < 85\,000$.
\shortans{70000}
\loigiai{
\textbf{Đáp số:} $70000$\\
$70000$ (hoặc $80000$).
}
\end{ex}

\begin{ex} %[3D11V60]
Có bao nhiêu số tự nhiên nằm giữa hai số $49\,998$ và $50\,002$?
\shortans{3}
\loigiai{
\textbf{Đáp số:} $3$\\
Các số là $49999, 50000, 50001$ (3 số).
}
\end{ex}

\begin{ex} %[3D11C60]
Tìm số tròn nghìn có 5 chữ số lớn nhất có tổng các chữ số bằng 12.
\shortans{93000}
\loigiai{
\textbf{Đáp số:} $93000$\\
Số có dạng $\overline{ab000}$. Để lớn nhất ta chọn $a = 9 \Rightarrow b = 12 - 9 = 3$. Số đó là $93000$.
}
\end{ex}
